\chapter{The condition of the planet and conditional futures}\label{ch:planet}
\question{What do environmental observations establish, and what do future pathways mean?}

An environmental statement needs a time, a place and a measurement. ``Warmer'' requires a reference period. ``Threatened'' is not the same as extinct. A projection describes what follows under specified assumptions; it does not observe the future. Keeping these distinctions clear makes the motivation stronger, because the reader can see what the evidence actually supports.

The examples here are selected assessment findings, not a complete planetary dashboard. They retain their publication dates and data periods. This edition does not present a historical assessment as a live September 2026 estimate.

\section{An observed change}
The IPCC's 2023 synthesis reports global surface temperature in 2011--2020 at 1.09 degrees Celsius above 1850--1900, with a very likely range of 0.95--1.20 degrees Celsius. The bracket describes the assessment's 5--95\% range. The same assessment attributes global warming unequivocally to human activities, principally greenhouse-gas emissions.\cite{ipcc}

This is a change in a global average over a decade. It is not the temperature in the clinic, a maximum temperature during a heatwave or a count of people exposed. An annual observation and a multi-year warming measure answer different questions. Changing the reference period without saying so would change the meaning of the number.

A global average is useful for describing a planetary change. A local decision requires further information. The clinic's water supply, building design and ability to operate during heat cannot be read from that average alone. This is an example of why a broad indicator and a local physical account can both be needed.

\section{Conditional pathways}
The IPCC also assesses contrasting greenhouse-gas scenarios. Its estimates for average warming in 2081--2100, relative to 1850--1900, include the following.\cite{ipcc}

\begin{center}\small
\begin{tabular}{@{}p{157pt}rr@{}}
\toprule
Scenario & Best estimate & Very likely range\\
\midrule
Very low emissions (SSP1-1.9) & $1.4\,^{\circ}\mathrm{C}$ & $1.0$--$1.8\,^{\circ}\mathrm{C}$\\
Intermediate (SSP2-4.5) & $2.7\,^{\circ}\mathrm{C}$ & $2.1$--$3.5\,^{\circ}\mathrm{C}$\\
Very high emissions (SSP5-8.5) & $4.4\,^{\circ}\mathrm{C}$ & $3.3$--$5.7\,^{\circ}\mathrm{C}$\\
\bottomrule
\end{tabular}
\end{center}

These are conditional climate responses to named pathways, not probabilities that society will choose each pathway. They are not a current-policy forecast. The ranges within a row describe uncertainty in the assessed response; the differences between rows also reflect different emissions assumptions. The IPCC assesses escalating risks with additional warming, with very high confidence.\cite{ipcc}

A reader should therefore resist two shortcuts. One is treating an adverse pathway as an inevitable destiny. The other is treating the existence of a lower-emissions pathway as proof that it will happen. Choices and constraints still matter. A table can clarify their consequences without making them for us.

\section{Biodiversity, materials and pollution}
In its 2019 global assessment, IPBES estimated that around one million animal and plant species were threatened with extinction. This is an assessment of threat, not a count of species already extinct. The estimate draws on assessed groups and estimated total diversity; insect threat is an important uncertainty. IPBES labels the supporting estimate established but incomplete.\cite{ipbes}

A threatened-species estimate has different units and meaning from a temperature change. It cannot be added to degrees Celsius to obtain a scientific ``damage total''. Habitat, species and ecosystem function also need more than a single worldwide count if we want to understand a particular place.

The UNEP International Resource Panel's 2024 outlook reports growth in material extraction and describes a pathway in which global extraction could rise about 60\% from 2020 to 2060 without a change of course. This is a conditional scenario, not a measured 2060 quantity. Its official overview supplies no uncertainty interval for that headline; none is invented here.\cite{resources}

Material quantity is itself incomplete as an impact measure. A tonne tells us mass. To judge consequences, we would also need to know what material, obtained where, by which process, for what use and with what residuals. The account must not acquire ecological meaning merely because it uses physical units.

Pollution adds another dimension. WHO's October 2024 assessment page describes a substantial worldwide health burden from outdoor air pollution using estimates for 2019.\cite{air} This book uses that qualitative finding without treating unlike health, climate and biodiversity measures as independent quantities that can simply be summed.

\section{Why this evidence matters to the question}
These findings make physical conditions difficult to dismiss as something outside an economy. Production, infrastructure and service depend on such conditions and can change them. Their effects are distributed unevenly, as the household example in Chapter~\ref{ch:access} illustrates conceptually.

Yet none of the assessments evaluates the EBU proposal developed here. They do not select its reference values, establish its scalar weights or test its actor policy. We should carry forward the motivation without borrowing authority for an untested remedy.

\practice{Is the difference between two scenario rows an observed temperature reduction caused by EBU?}{No. Both rows describe conditional future climate assessments. Neither includes an EBU intervention. Their contrast explains why assumptions matter, not what this project has achieved.}

\carry{Diagnosis gives us reasons to ask for better physical visibility. The next step is to specify what a proposed account would need to contribute and how it could be found inadequate.}
